\chapter{Discussion}
\label{ch:discussion}

\Cref{ch:results} reported what was observed. This chapter asks what it means. It works through
the four research questions in turn (\Cref{sec:disc-capacity,sec:disc-policy,%
sec:disc-bottleneck,sec:disc-modes}), relates the findings to the literature
(\Cref{sec:disc-literature}), draws out what follows for practice
(\Cref{sec:disc-managerial}), and closes with what the study cannot support
(\Cref{sec:external-validity,sec:disc-limitations}).

\section{What simulation adds to capacity planning (RQ1)}
\label{sec:disc-capacity}

The most consistent result of this study is that the analytical capacity estimate exceeded the
simulated recommendation in every one of the twelve months, by as much as nine full-time
equivalents in December. That direction is worth dwelling on, because it is the opposite of
what one might expect.

The intuitive argument for simulating rather than calculating is that a formula \emph{under}
-estimates requirements: it ignores queueing, and queueing makes things worse. That argument is
sound as far as it goes, and it is why the analytical estimate here targets \SI{85}{\percent}
utilisation rather than \SI{100}{\percent} --- a margin motivated precisely by the
utilisation--congestion relationship that makes waiting time grow without bound as utilisation
approaches one \citep{hopp2011}. The safety margin accounts for part of the observed gap.

It does not account for all of it, and the remainder is the more interesting finding. The
analytical estimate sizes capacity as though every order were interchangeable. It has no
representation of the fact that a sequencing policy can decide \emph{which} orders wait. Once
the evaluation allows an urgency-aware discipline, a smaller workforce becomes sufficient,
because the scarce capacity is directed at the orders whose deadlines are binding. What the
simulation discovers, in other words, is a \emph{substitution between capacity and operating
discipline} — and that substitution is invisible to any formula that reasons about aggregate
workload.

This reframes the contribution of simulation to capacity planning. It is not merely a more
accurate calculator of the same quantity. It evaluates a fundamentally larger decision space,
one in which how work is ordered is a decision variable rather than a fixed background
condition. A planner using the formula alone would be choosing among workforce sizes; a planner
using the framework is choosing among (workforce, discipline) pairs, and the best pair is
cheaper than the best size.

The practical magnitude is not trivial. Nine FTE in December is roughly \SI{19}{\percent} of
the analytical estimate of 47, and even the modest one-FTE gaps in the summer months are real
money. But the direction carries a caution as well: because the framework
recommends leaner configurations than the formula, it depends more heavily on the sequencing
policy actually being implemented as modelled. A recommendation of 38 FTE operating under RL-3
is not a recommendation of 38 FTE operating under whatever discipline the floor happens to use.

\section{Sequencing policy and differentiated service (RQ2)}
\label{sec:disc-policy}

\subsection{Policy value is conditional on the capacity regime}

The headline finding on sequencing is not that one policy dominates. It is that the
\emph{magnitude} of the difference between policies depends almost entirely on how tightly
capacity is constrained.

At October's \rgm{s11\_7\_5} configuration, all three policies exceeded \SI{99.9}{\percent}
attainment on both classes and their total costs differed by \eur{30} --- five hundredths of
one per cent. At October's \rgm{s753} configuration, in the same month with the same demand,
FIFO delivered \SI{24.5}{\percent} urgent attainment while the alternatives delivered essentially
\SI{100}{\percent}. The system is the same; only the capacity is different.

The mechanism is straightforward once stated. A sequencing policy can only act when there is a
queue to choose from. When capacity is ample, queues are short and transient, decision points
are rare, and every policy makes broadly the same choices because there is usually only one
order waiting. When capacity is tight, queues are long and persistent, decision points are
continuous, and the policy determines the entire allocation of waiting across the order
population.

The consequence for practice is that ``which sequencing policy should we use?'' is the wrong
question asked in isolation. The right question is whether the operation runs close enough to
its capacity limit for the answer to matter --- and, in this study's demand profile, it does so
in roughly a third of the year.

\subsection{Why FIFO fails, and why capacity did not rescue it in the evaluated range}

FIFO's failure in the four campaign months is categorical rather than marginal: it satisfied
the service floors at none of the sixteen candidate configurations in January, February,
November or December. Elsewhere in the year its failures are partial --- \num{44} infeasible
configurations spread over the eight remaining months --- so what distinguishes the campaign
months is
that no candidate at all worked. This deserves careful interpretation, because it is tempting to
read it as ``FIFO simply needs more capacity''.

That reading is not supported by the range examined. The candidate sets in those months extend
well above the analytical estimate, and FIFO fails throughout. The reason is structural. Under FIFO an urgent
order entering the queue behind a large volume of normal orders must wait for all of them, and
its deadline is six times tighter (\SI{240}{\minute} against \SI{1440}{\minute}). Because
urgent orders are a minority of volume even in December --- roughly a quarter --- the expected
wait ahead of any urgent order is dominated by normal orders. Adding capacity shortens all
queues proportionally; it does not change the fact that the urgent order's position in the
queue bears no relation to its deadline. To satisfy a \SI{240}{\minute} threshold under FIFO,
one would need enough capacity to keep essentially \emph{every} queue short. No configuration
in the evaluated range came close, and the capacity that would be required lies far outside the
economically sensible band.

The claim this supports is bounded, and the boundary matters. FIFO is not incapable of meeting
differentiated deadlines in principle --- with a large enough workforce, queues shorten until
even the least favourably positioned urgent order clears its threshold. What the experiment
establishes is that within the sixteen workforce candidates evaluated, in each of the four
campaign months, no such configuration existed, so the shortfall could not be closed by the
capacity levers a planner would actually consider. This is a concrete instance of a point well
established in scheduling theory but easy to lose in practice: when deadlines are
differentiated, a discipline that ignores the differentiation is expensive to repair by
resourcing alone \citep{panwalkar1977}. It is cheaper to encode the commitment in the operating
rule.

\subsection{What the learned policy actually contributes}

RL-3 was the lowest-cost policy in seven of twelve months and Urgent-First in five, which on
its own would be an unremarkable split. The substantive finding is narrower and more specific.

Both urgency-aware policies protect the urgent class almost completely --- the class-level
trade-off plot shows them compressed into a narrow band near \SI{100}{\percent} urgent
attainment. They differ in what that protection costs the normal class. Strict priority, by
construction, gives the normal class whatever capacity remains after every urgent order has been
served. When the urgent share is small this is inexpensive; when it grows, the normal class
absorbs the entire shortfall.

December is where this becomes decisive. At the recommended configuration, Urgent-First achieved
\SI{99.3}{\percent} urgent attainment but drove normal attainment to \SI{73.8}{\percent},
\emph{below the \SI{80}{\percent} floor}, and was therefore infeasible. RL-3 achieved
essentially identical urgent attainment (\SI{99.2}{\percent}) while holding the normal class at
\SI{85.4}{\percent}. It was the only feasible policy at that configuration, and it was
\SI{13.2}{\percent} cheaper.

What the learned policy contributes, then, is not better urgent-class protection --- strict
priority already achieves that. It is the ability to \emph{stop} protecting the urgent class at
the point where doing so costs more than it gains. The stage-differentiated decision shares
observed in the forecast run are consistent with this: the agent chose the urgent queue in
\SI{26.7}{\percent} of picking decisions but \SI{100}{\percent} of dispatch decisions. Picking
is where the queue is long and where deferring an urgent order is cheap in slack terms; dispatch
is the last opportunity to rescue an order's deadline. A fixed rule cannot make that
distinction, because it applies the same logic everywhere.

This should not be over-read. The agent was not designed to discover that behaviour, and this
study cannot establish that it represents a general principle rather than an artefact of these
training conditions. What can be said is that a policy conditioned on stage, queue state and
remaining slack found a balance that a stage-independent rule could not, in the specific regime
where the binding constraint moved from the urgent class to the normal one.

\subsection{Where the learned policy fails, and why it matters less than it appears}
\label{sec:disc-rl-failure}

The systematic comparison in \Cref{sec:res-matched} qualifies the above in two ways that a
reader inclined to favour the learned policy should take seriously.

The first is that RL-3 achieved higher urgent-class attainment than strict priority in
\emph{none} of the 192 matched configurations. This is not a defect --- it is what a strict
priority rule guarantees by construction --- but it does bound the claim that can be made. RL-3
does not protect the urgent class better; it protects it very slightly less well, by
\num{0.62} percentage points on average, and trades that for normal-class service.

The second is more substantial. Averaged over all 192 configurations, RL-3's normal-class
attainment is \emph{worse} than Urgent-First's, because in 25 configurations it loses heavily
--- by up to 34 percentage points. Reported without qualification, that would undermine the
entire argument of \Cref{sec:res-policy}.

The qualification is that all 25 of those configurations are ones in which \emph{neither} policy
produced a feasible plan. They are severely under-resourced: backlog between 27 and 35 per cent
of the month's orders, picking utilisation averaging \SI{93.8}{\percent}. In May at
\rgm{s232}, for instance, the three policies deliver urgent attainment of \SI{4.0}{\percent},
\SI{100}{\percent} and \SI{94.1}{\percent} and normal attainment of \SI{27.3}{\percent},
\SI{57.4}{\percent} and \SI{23.7}{\percent}. No manager would deploy any of them. In the 167
configurations where at least one policy produced a viable plan, RL-3's worst normal-class
result was \num{0.08} percentage points below Urgent-First --- effectively never worse --- and
its average was \num{1.84} points better.

Two readings of this are available and they should be distinguished. The favourable reading is
that RL-3's weakness is confined to a region of the decision space that is irrelevant to
practice, so it does not detract from its usefulness where plans are actually made. The
unfavourable reading is that a learned policy degrading sharply outside the conditions where it
performs well is precisely the reliability concern raised by \citet{waubert2022}: the agent was
trained across a range of workforce configurations, and its behaviour deteriorates at the edge
of that range in a way a two-line heuristic's does not.

Both readings are defensible on this evidence, and this study cannot adjudicate between them,
because it did not test the agent under conditions systematically outside its training
distribution. What can be said is that the failure mode is identified, bounded and explicable,
which is more than an aggregate performance figure would have revealed. It also suggests a
concrete safeguard for any deployment: a learned sequencing policy should be paired with a
feasibility check on the configuration it is running under, and should not be trusted in regimes
where that check fails.

\subsection{The Urgent-First implementation and its within-class ordering}
\label{sec:disc-uf-ordering}

One implementation detail bears on how the Urgent-First results should be read. As established
in \Cref{sec:urgent-first}, the priority queue orders orders strictly by class but leaves the
sequence \emph{within} a class unspecified --- it follows the internal structure of a binary
heap, which is neither first-in-first-out nor sorted by arrival.

This is not a neutral detail for the normal class. Under a first-in-first-out discipline within
the class, normal orders would wait in a predictable order and their waiting times would be
comparatively concentrated. Under heap ordering, a normal order can be overtaken repeatedly by
later arrivals of its own class, which lengthens the \emph{tail} of the waiting distribution
without necessarily increasing its mean.

The observed waiting times show exactly that signature. At January's recommended configuration
\rgm{s13\_8\_4}, picking utilisation is identical to three decimal places across all three
policies (\num{0.958}), so any difference is attributable to ordering rather than to capacity.
Under Urgent-First the \emph{mean} picking wait is \SI{301.8}{\minute} --- \emph{lower} than
FIFO's \SI{380.4}{\minute} --- while its p95 wait is \SI{1796.9}{\minute}, more than double
FIFO's \SI{842.4}{\minute}. RL-3, whose sub-queues are first-in-first-out, sits between them at
a p95 of \SI{993.8}{\minute}. Urgent-First is therefore serving the average normal order sooner
than FIFO while leaving a minority of normal orders waiting far longer than either alternative.

Because the service-level criterion is a threshold, the tail is what matters. Part of
Urgent-First's normal-class shortfall is therefore attributable to within-class ordering rather
than to the priority rule itself. A variant implementing first-in-first-out within each class
would plausibly narrow the gap to RL-3 --- and the fact that RL-3's engine \emph{does} maintain
first-in-first-out sub-queues means the comparison is not perfectly clean on this dimension.
This is stated as a limitation rather than corrected, in keeping with the decision to document
the system as implemented. It is the single most valuable target for future work identified by
this study.

\section{Explainability and the marginal FTE (RQ3)}
\label{sec:disc-bottleneck}

\subsection{One stage dominates}

Picking was the primary pressure location in all twelve months, and the concentration is
extreme: between \SI{83.7}{\percent} and \SI{100}{\percent} of the waiting time accumulated by
late orders occurred there. Packing and dispatch contribute almost nothing to lateness, despite
often running at substantial utilisation.

This is a direct consequence of the workload structure documented in
\Cref{sec:order-heterogeneity}. Picking carries roughly three times the workload units of
dispatch per order, and seasonal growth falls disproportionately on it --- picking workload per
order rose \SI{104}{\percent} between May and December against dispatch's \SI{7}{\percent}. A
model that assumed equal workload across stages would have distributed pressure evenly and
produced a materially different, and wrong, diagnosis.

That the diagnostic identifies one stage so consistently is convenient but should not be taken
as evidence that the diagnostic is sensitive. A weighted score in which one stage dominates
every component would rank that stage first under most plausible weightings. The Stage Pressure
weights were chosen for interpretability, not calibrated, and this study does not test whether
they would discriminate correctly in a more balanced operation. The appropriate claim is that
the diagnostic identified where pressure was concentrated in these scenarios, not that it is a
validated instrument.

\subsection{Diagnosis is not prescription}

The most managerially useful result of this study is arguably the negative one. In nine months
the adaptive search tested adding one worker at the identified bottleneck, and in all nine it
rejected the addition.

The gap is usually large. An FTE costs \eur{2400} per month; in six of the nine months it
avoided less than \eur{130} of late-order penalties. December is the clearest case: picking runs
at \SI{90.5}{\percent} utilisation and accounts for \SI{83.7}{\percent} of late orders' waiting
--- an unambiguous bottleneck by any reading --- and yet the tenth picking worker prevented
\eur{20} of penalties.

The reason is that the recommended configuration has already absorbed the failures that
capacity can cheaply prevent. What remains are failures concentrated in demand bursts that a
marginal worker cannot smooth, and orders whose deadlines were never going to be met. The
marginal return on capacity at that point is close to flat, while its cost is linear.

This has a clear implication for practice: \emph{a bottleneck is a description of where pressure
is, not an instruction to add capacity there}. The two are routinely conflated, and the
conflation is expensive. The value of pairing the diagnostic with an explicit economic test is
that it separates them. The framework tells a planner where the constraint is, and then tells
them --- separately, and usually in the negative --- whether relieving it is worth paying for.

November is the instructive exception, recovering \eur{2025} against a cost of \eur{2400}. At
that margin the decision is genuinely close, and factors outside the model --- risk tolerance,
the reputational cost of missing a peak, the option value of trained capacity going into
December --- could reasonably tip it. Surfacing that a decision is close is itself useful
output.

\section{Two workflows, two kinds of answer (RQ4)}
\label{sec:disc-modes}

The retrospective and forecast workflows share an engine, a feasibility criterion and an
economic model, but they do not produce the same kind of statement, and treating their outputs
as interchangeable would be a mistake.

The retrospective workflow holds demand fixed and produces a point answer: given this demand,
this configuration would have been advised. Its uncertainty structure is essentially absent ---
one realisation per month --- which is acceptable for its purpose, because the question is
counterfactual rather than predictive. Common random numbers make the \emph{comparison between
policies} fair at a given configuration, but they do not tell us how much the recommendation
would move under a different realisation of demand. That is a genuine limitation of the
retrospective evidence in this study.

The forecast workflow treats demand as uncertain and produces a spread-aware answer: this
configuration has this mean cost across replications, cost as much as that in the least
favourable of them, and met the service floors in this share of them. The two-stage screening
and validation design is a pragmatic response to cost --- evaluating every candidate on every
replication would be wasteful when most candidates are eliminated on the first --- and it makes
the consistency of the recommendation explicit rather than implicit.

The strength of that statement should not be overstated, and the constraint is the replication
count rather than the design. Three replications support a description of the spread observed;
they do not support an estimate of a cost distribution's tail, and the share of replications
meeting the floors is a success count over three trials rather than a probability with any
precision attached. The forecast workflow as exercised here therefore demonstrates that the
framework can express scenario variation and act on it, which is what RQ4 asks; it does not
deliver calibrated reliability figures, and \Cref{sec:future-work} identifies the replication
budget and genuine forecast-error propagation as the two extensions that would.

That the December forecast run reproduced the retrospective run's qualitative policy pattern
under an independent demand realisation and a different economic parameter set is mild evidence
that the pattern is a property of the system rather than of one realisation. It is only mild
evidence: both runs share the same generator, the same service-time model and the same
structural assumptions, so agreement between them is not independent corroboration.

A final point of discipline. The two runs used different economic parameters, and no absolute
monetary comparison between them appears anywhere in this thesis. This is not pedantry: the
retrospective run penalises a late urgent order at \eur{20} and the forecast run at \eur{15},
while valuing labour at \eur{15} and \eur{18} per hour respectively. Costs from the two runs are
denominated in different assumptions and are simply not the same quantity.

\section{Relation to the literature}
\label{sec:disc-literature}

The findings connect to three strands of the literature, and the connections are not uniformly
confirmatory.

\textbf{Warehouse operations research} has concentrated overwhelmingly on the internal
efficiency of picking --- storage assignment, routing, batching and zoning
\citep{dekoster2007,gu2007,vangils2018}. That focus is well justified by the cost share of
picking, and this study's finding that picking dominates pressure in every month is consistent
with it. But the response studied here is different. Where that literature asks how to make
picking faster, this study takes the picking process as given and asks how much of it to buy
and in what order to feed it. The two are complementary: an improvement in picking productivity
would shift the capacity requirement without changing the structure of the sequencing decision.

\textbf{The call for integration} across planning problems is directly relevant.
\citet{vangils2018} argue that order-picking planning problems are usually studied in isolation
and that combining them yields better outcomes than optimising each separately, a claim
reinforced by their later analysis of real-life features \citep{vangils2019}. This study
provides evidence for an analogous claim one level up: capacity and sequencing, treated
jointly, produce a cheaper feasible plan than capacity treated alone --- which is precisely what
the systematic gap between the analytical estimate and the simulated recommendation measures.

\textbf{Scheduling and priority rules.} That priority dispatching outperforms
first-in-first-out when due dates are differentiated is long established
\citep{panwalkar1977}, and the results here are consistent with it. The finding that is
\emph{not} simply a restatement concerns the boundary of that advantage: strict priority itself
fails when protecting the priority class pushes the other class below its own floor. The
classical literature largely evaluates rules against aggregate objectives such as mean tardiness
or makespan, under which a strict priority rule looks uniformly good. Under a conjunctive
class-specific constraint it does not, and December demonstrates the failure concretely.

\textbf{What queueing theory predicted, and the model confirmed.} Two results in this study are
best read as empirical confirmations that the model behaves as theory requires, rather than as
discoveries. The first is that backlog and utilisation are almost invariant to sequencing policy
(\Cref{sec:res-backlog}). Capacity fixes the worker-minutes available in the month, and no
ordering of the queue creates more of them; the conservation properties of non-pre-emptive
priority disciplines and Little's law \citep{little1961} both express that constraint. They do
not, however, guarantee that the \emph{number} of orders finishing before a finite cut-off is
identical under every discipline: service times are heterogeneous here, so a discipline that
happens to favour shorter jobs can complete slightly more of them before the horizon. The
observed figures show exactly that residual --- backlog shares of \SI{8.59}{\percent},
\SI{8.97}{\percent} and \SI{8.95}{\percent} in January across the three policies differ, but by
fractions of a percentage point. The empirical claim, which is the one this thesis makes, is
that at the configurations tested the sequencing policies produced almost identical utilisation
and backlog, so their operative effect was the allocation of lateness between classes rather
than any material change in aggregate completion. Observing that invariance is a useful internal
validation check --- had it failed by a wide margin, something in the engine would have been
wrong.

The second is the shape of the utilisation--congestion relationship. The measured \emph{absolute}
cost difference between the two urgency-aware policies rises monotonically across the five
observed utilisation bands, by a factor of nearly twelve from below \SI{70}{\percent} to above
\SI{95}{\percent}, and by a factor of roughly ten once normalised per \num{1000} orders
(\Cref{tab:pressure-bands}). This is consistent with the non-linear growth of waiting with
utilisation \citep{hopp2011}: sequencing can only act on queues, queues barely form at low
utilisation, and their length grows sharply as utilisation approaches one. Two cautions belong
with it. The measured quantity is divergence between the policies, not RL-3's advantage; within
the feasible region the signed difference does favour RL-3 and does grow with pressure, but
where no policy reaches feasibility it reverses. And utilisation is a grouping variable here,
not a demonstrated cause --- the high bands coincide with campaign months, high urgent shares
and lean configurations, which this design cannot separate. The contribution is therefore not
the mechanism, which is textbook, but its \emph{quantification} in a specific decision context:
a planner can read from \Cref{tab:pressure-bands} roughly how much separation between the two
policies to expect at their own utilisation, and whether their operation sits in the region
where that separation favours the learned policy.

\textbf{Workload control.} The closest conceptual relative to this work outside the warehousing
literature is workload control, which manages congestion by regulating the release of work into
a shop rather than by expanding capacity \citep{thurer2012}. Both approaches treat the timing and
ordering of work as a lever operating on a fixed capacity base, and both are motivated by the
observation that this lever is cheaper than capacity. The results here support that motivation
quantitatively: in the four campaign months, changing the sequencing rule was the difference
between no viable plan at any of sixteen workforces and a viable plan at almost all of them ---
an outcome none of the capacity increases within that candidate range achieved.

The differences are equally instructive. Workload control decides \emph{when} work is released;
the policies here decide \emph{which} waiting work is served next, and release is exogenous. And
where workload control typically takes capacity as given, this framework treats capacity as the
co-equal second decision. A natural synthesis --- controlling release, sequencing and capacity
jointly --- is beyond what was studied here but follows directly from putting the two literatures
side by side.

\textbf{Learning-based control.} The use of a deep Q-network to select among sequencing
decisions places this work alongside a growing literature applying reinforcement learning to
dynamic scheduling \citep{luo2020,liu2024}. The comparison should be made carefully. That
literature typically pursues policy performance as the object of study, on benchmark instances,
against objectives such as tardiness. Here RL-3 is one of three compared components inside a
capacity-planning framework, and the interesting result is not that it performs well but
\emph{where} it does: in a narrow regime where the binding constraint shifts from one class to
the other. In most of the year it is indistinguishable from a rule that takes one line to
describe. Any assessment of whether the additional complexity is warranted must weigh that
narrowness against the cost of training, maintaining and explaining a learned policy --- and in
a setting where a manager must defend the plan, explainability is not a minor consideration.

\section{Managerial implications}
\label{sec:disc-managerial}

Five implications follow for someone responsible for a monthly fulfilment plan.

\textbf{Size capacity and discipline together.} The workforce required depends on how work is
ordered. Sizing headcount from workload ratios and then choosing an operating rule separately
over-provides capacity --- here by up to \SI{19}{\percent} of the estimate in the peak month.

\textbf{Never accept an aggregate service figure.} A configuration reporting \SI{81}{\percent}
overall attainment delivered \SI{24.5}{\percent} to its urgent customers. Where commitments are
differentiated, they must be measured and constrained separately.

\textbf{Know whether you are in the regime where sequencing matters.} If the operation runs
comfortably below capacity, the operating rule is close to irrelevant and effort is better spent
elsewhere. If it runs near capacity in peak months --- as most seasonal fulfilment operations do
--- the rule determines whether the commitments can be met at all.

\textbf{Do not treat a bottleneck as a hiring case.} The constraining stage is where to look,
not automatically where to spend. In nine of nine tests here, the additional worker at the
identified bottleneck did not pay for itself. Ask what it would prevent, and price that.

\textbf{Plan peaks across scenarios, not on a single one.} For a forecast month, a mean cost is
less useful than the share of plausible demand realisations in which the plan holds. A
configuration that succeeds on average but fails in a third of the scenarios tested is not a
plan. Establishing that share with any precision needs considerably more than the three
replications used here.

\subsection{Using the pressure bands as a screening test}

One result is directly portable to an operation that has not built this framework.
\Cref{tab:pressure-bands} relates how far apart the two urgency-aware policies land to a
quantity most warehouses already measure: utilisation at the constraining stage. Below
\SI{70}{\percent} the mean absolute difference between them was \eur{537} per month; above
\SI{95}{\percent} it was \eur{6320}. The association is observed, not established as causal,
and the direction of the difference depends on whether a feasible plan exists at all.

A manager can therefore use their own utilisation figure as a cheap screening test for whether
this kind of analysis is worth commissioning at all. If the constraining stage runs comfortably
below \SI{80}{\percent} through the year, the evidence here suggests the sequencing rule is
close to a free choice and effort is better directed elsewhere --- at demand shaping, at
process improvement, or at the storage and routing decisions that the order-picking literature
addresses \citep{petersen2004}. If it runs above \SI{90}{\percent} in the peak months, as
seasonal fulfilment operations typically do, the operating discipline is likely to be worth more
than its analysis costs.

The threshold is specific to this model and this demand profile, and should be treated as an
order-of-magnitude guide rather than a calibrated rule. What is transferable is the shape of the
relationship, which follows from queueing behaviour rather than from anything peculiar to the
implementation.

\subsection{What the framework does not tell a manager}

It is worth being explicit about the questions a planner might reasonably bring to this
framework and not get an answer to, since a decision-support tool that is unclear about its own
boundaries invites misuse.

It does not say whether the recommended capacity can be \emph{staffed}. A recommendation of 22
picking FTE says nothing about whether 22 months of picking labour can be arranged into a legal,
workable shift pattern given the operation's contracts, absence rates and recruitment lead
times. That translation is a separate problem, and a configuration this framework considers
cheapest may be one the operation cannot actually roster.

It does not say whether the service thresholds are the right ones. The framework takes
\SI{240}{\minute} and \SI{1440}{\minute}, and the \SI{95}{\percent} and \SI{80}{\percent}
attainment floors, as given. Whether those commitments are commercially sensible --- whether the
urgent tier is priced to cover the capacity it requires --- is a question the model can inform
by costing alternatives, but cannot answer.

It does not say what happens outside the month. The finite horizon makes backlog visible but
does not carry it forward: December's unfinished orders do not become January's arrivals.
For an operation with sustained overload across consecutive months, that omission would
understate the problem.

And it does not detect a wrong demand forecast. The forecast workflow evaluates robustness
across realisations \emph{consistent with the volume assumption supplied to it}. If that
assumption is wrong, every candidate is evaluated against the wrong month and the reported
reliability figures are correspondingly misleading.

\section{External validity}
\label{sec:external-validity}

The framework as it stands is an academic prototype, and the distance between it and
operational deployment should not be understated.

Deployment would require, at minimum: service-time distributions fitted to observed handling
data rather than assumed; economic parameters calibrated to the operation's actual labour costs
and contractual penalties; service thresholds expressed against the operation's real calendar,
shift pattern and carrier cut-offs rather than an abstracted operating clock; a demand forecast
with a validated error structure, since the forecast workflow's output is only as good as the
volume assumption entering it; and a translation layer converting monthly FTE recommendations
into feasible rosters, which may itself be infeasible for a configuration the model considers
optimal.

Beyond calibration, external validation would require comparing the model's recommendations
against outcomes in a real operation over multiple planning cycles --- something no part of this
study attempts. Until that is done, the appropriate claim is that the framework is a coherent
instrument for comparing configurations under stated assumptions, not that it predicts the
performance of any particular warehouse.

\section{Limitations}
\label{sec:disc-limitations}

The methodological limitations were stated in \Cref{sec:method-limitations}. Four bear
repeating because they bound the interpretation of the results specifically.

\textbf{The evidence is synthetic.} Every quantitative finding is a property of a model
operating on generated data. The patterns are internally consistent and mechanistically
explicable, but none is an empirical observation about a real warehouse.

\textbf{The two engines differ in within-class mechanics.} RL-3 maintains first-in-first-out
sub-queues; the Urgent-First baseline does not. Part of the observed gap between them in
normal-class attainment is plausibly attributable to that difference rather than to the decision
rule, and this study cannot separate the two effects.

\textbf{Retrospective results rest on one realisation per month.} No confidence interval
accompanies any retrospective recommendation. A different demand realisation could shift a
marginal month's recommended configuration, and the near-ties observed in the shoulder months
--- August differed by \eur{10} --- are well within the range where this matters.

\textbf{The learned policy cannot be regenerated.} The trained network weights are not retained
in the repository. The persisted evaluation outputs are available and every RL-3 result reported
here derives from them, but exact regeneration requires the checkpoint, which was not available
in the environment used for final thesis preparation; retraining from the recorded configuration
would produce a different policy rather than the same one. The results are internally consistent
and were generated by the documented pipeline, but no claim of end-to-end reproducibility of the
RL-3 results from the distributed artefacts is made anywhere in this thesis. This is a genuine
weakness affecting one of the three compared policies, and it is recorded as such rather than
minimised.

\textbf{The state representation saturates at peak.} The learned policy's capacity and
queue-length features are clipped, and at December's candidate workforces the picking-capacity
feature and, during congestion, the picking queue-length feature both sit at their ceiling
(\Cref{sec:app-saturation}). The agent therefore had reduced resolution on exactly the
configurations where the sequencing decision mattered most. It nevertheless produced feasible
lower-cost outcomes there, so this bounds the representation rather than invalidating the
result, but it means the reported behaviour should not be extrapolated to capacities well beyond
the normalising constants.
